% -----------------------------------------------------------------------------
% Introduction
% -----------------------------------------------------------------------------
\section{Introduction}
\label{sec:introduction}
\PARstart{I}{nsider} threats are hard to detect because authorised users
can misuse legitimate access without immediately standing out from normal
organisational activity. Such misuse may expose sensitive information,
disrupt operations, or weaken organisational trust. It can also build up
gradually through small changes in logon, file, communication, web, or
removable-media behaviour that look harmless on their own
\cite{AlzaabiMehmood2024,AlketbiMehmood2025}. The significance of these
changes often becomes visible only once related actions are examined
together over time \cite{Eshmawi2026,Pal2023}.

\smallskip
Insider-threat activity records are heterogeneous, high-dimensional, and
severely imbalanced: normal behaviour accounts for almost all observations,
and malicious activity forms only a small minority
\cite{AlShehariAlsowail2023,YuanWu2021}. Millions of events must be
transformed into manageable behavioural representations without discarding
inactive periods, chronological order, or gradual transitions
\cite{LeZincirHeywood2021,Alshehri2022,Xiao2024}. Accuracy alone can hide
poor rare-class detection, so evaluation should emphasise PR-AUC, precision,
recall, F1-score, false positives, false negatives, and the resulting alert
burden.

\smallskip
Random Forest (RF) and Extreme Gradient Boosting (XGBoost) remain strong
references once behaviour has been converted into engineered tabular
features. Grinsztajn, Oyallon, and Varoquaux showed that tree ensembles
remain competitive with deep neural models on typical tabular datasets
\cite{Grinsztajn2022}. Their benchmark did not cover the stream-like,
chronological, and extremely imbalanced conditions studied here, so it
supports including demanding tree baselines without predicting their
performance on chronological insider-threat data.

\smallskip
Bidirectional long short-term memory (Bi-LSTM) networks can capture
relationships across ordered observations, and attention mechanisms
aggregate evidence across a behavioural window \cite{YuanWu2021,Song2024}.
Existing security studies often combine representation learning,
augmentation, temporal analysis, and classification through separately
trained or sequential stages \cite{He2024,Gayathri2024}. Such modular
pipelines join complementary strengths, but they do not provide a single
gradient pathway from the final classification objective through both the
temporal representation and the downstream decision mechanism. A
differentiable classifier head permits that shared pathway, which raises a
more specific question: can the temporal encoder and the differentiable
decision mechanism be updated jointly without destabilising useful
prediction and routing behaviour?

\smallskip
Teacher--student learning offers another route to stable optimisation. In
knowledge distillation, a trained teacher supplies prediction information
that complements supervision from the true labels
\cite{Hinton2015Distillation,Gou2021KDSurvey}, and the student need not be
smaller than the teacher \cite{Furlanello2018BornAgain}. Security studies
have used teacher guidance for privacy-preserving knowledge transfer,
lightweight intrusion detection, and anomalous user-activity analysis
\cite{Shafee2020Mimic,Hsu2023TeacherStudent,Umair2025KDIDS}. Their principal
purposes were prediction transfer, multi-label learning, privacy, or
compression. The reviewed studies do not directly evaluate the preservation
of differentiable tree-routing behaviour while a same-architecture temporal
student is optimised jointly.

\smallskip
The reviewed literature does not establish a controlled procedure for
updating a temporal differentiable-tree model end-to-end while constraining
drift in its prediction logits and its ODST route probabilities relative to
a frozen same-architecture reference. It also leaves open whether the
resulting tree-contribution rankings remain meaningfully connected to the
model's prediction function under explicit intervention. This paper
therefore targets optimisation stability and intervention-tested explanation
faithfulness for rare, chronologically ordered insider-threat sequences,
rather than predictive superiority.

\smallskip
We use a frozen, same-architecture teacher to guide the joint optimisation
of a temporal sequence--ensemble student. The student combines a Bi-LSTM
encoder, temporal attention, and an oblivious differentiable sparsemax tree
(ODST) head. It is initialised from the teacher and trained on the true
labels together with logit- and route-consistency losses. The logit term
penalises deviation from the teacher's detached prediction logits, and the
route term penalises deviation from its detached ODST routing probabilities;
we call this procedure teacher-anchored logit-and-route regularisation. Its
aim is not to raise predictive performance beyond the teacher, but to let
every trainable component of the student update while its prediction-logit
and routing drift are constrained. The teacher is used only during training
and removed at inference, so it serves as an optimisation reference rather
than an inference component or a model-compression mechanism.

\smallskip
The study addresses three research questions:
\begin{enumerate}
    \item[\textbf{RQ1:}] Under the same input content, what predictive
    differences arise between sequence-specific and fixed-position
    flattened models, and how do chronological order and available
    behavioural history affect trained sequence models evaluated without
    retraining?
    \item[\textbf{RQ2:}] Can teacher-anchored logit-and-route
    regularisation permit reproducible end-to-end student updates while
    retaining teacher-level validation performance and limiting
    prediction-logit and ODST-routing drift across predefined random seeds?
    \item[\textbf{RQ3:}] What evidence does the resulting architecture
    provide concerning explanation faithfulness, probability calibration,
    alert burden, computational cost, and controlled degraded-input
    sensitivity?
\end{enumerate}

\smallskip
The experimental design separates development, reference-model testing,
teacher-anchored validation, sensitivity, and post-hoc evidence. Synthetic
Carnegie Mellon University Computer Emergency Response Team (CERT) r4.2
supports representation development, routing diagnosis, explanation
development, and controlled degraded-input analysis. A predefined one-pass
CERT r5.2 test evaluates the RF, XGBoost, attention--linear, and
frozen-encoder ODST reference families, whereas the final teacher-anchored
students are evaluated separately on the CERT r5.2 train/validation split.
Because that test partition had already been accessed for the reference
models, any later final-student analysis on it would constitute post-hoc
corroboration rather than independent test evidence. Full protocol details,
including the same-user partition structure and predefined viability
margins, are given in Section~\ref{sec:methodology}.

\smallskip
Beyond this core comparison, the paper reports four further analyses: a
same-input-content comparison of sequence-specific and flattened
representations; controlled reversal, shuffling, and partial-history
interventions that probe dependence on chronological order; a multi-day
malicious-activity cohort; and a reference-centred faithfulness evaluation
of ODST tree contributions. The corresponding predictive, intervention, and
sensitivity results appear in Section~\ref{sec:experimental_results}. Their
combined interpretation is component-specific: Bi-LSTM--attention provides
the principal temporal representation pathway; teacher anchoring provides
the stabilising optimisation procedure; ODST provides an
explanation-oriented differentiable head with traceable route, leaf, and
tree-level outputs; and attention--linear remains the efficiency- and
raw-calibration-oriented neural reference.

\smallskip
The paper makes four contributions:
\begin{itemize}
    \item A chronological behavioural representation that retains inactive
    days and prevents sliding windows from crossing temporal partition
    boundaries, complemented by reversal, shuffling, and partial-history
    interventions.
    \item A teacher-anchored logit-and-route regularisation procedure
    evaluated across three predefined CERT r5.2 validation seeds for
    predictive non-degradation and routing stability during end-to-end
    student optimisation.
    \item A same-information comparison in which sequence-aware and
    fixed-position flattened models receive identical input values, so the
    model families can be compared without changing input content.
    \item An intervention-based evaluation showing that ODST's
    reference-centred tree-contribution rankings are faithful under the
    defined top-versus-random replacement test; complementary analyses
    report route and leaf traceability, probability calibration, alert
    consolidation, latency, and controlled degraded-input sensitivity.
\end{itemize}